% ============================================================
% Capitulo 3 -- Metodologia
% Tamanho-alvo: 8-12 paginas
% Ordem de escrita: PRIMEIRO (junto com Resultados).
% Mais denso em fatos verificaveis = menos espaco pra hand-waving.
% ============================================================

\chapter{Metodologia}
\label{cap:metodologia}

% ------------------------------------------------------------
\section{Vis\~ao geral do pipeline}
\label{sec:pipeline}

% ESCREVER (1 paragrafo): descricao breve dos blocos do pipeline.
% A figura \ref{fig:pipeline} (F0) abaixo e o diagrama em TikZ.

\begin{figure}[H]\centering
\resizebox{\linewidth}{!}{%
\begin{tikzpicture}[
    node distance=0.9cm and 0.7cm,
    box/.style    ={rectangle, rounded corners, draw, thick,
                    minimum width=2.0cm, minimum height=0.9cm,
                    align=center, font=\scriptsize},
    data/.style   ={box, fill=blue!10},
    proc/.style   ={box, fill=green!15},
    model/.style  ={box, fill=orange!18},
    eval/.style   ={box, fill=red!12},
    arr/.style    ={->, >=stealth, thick}
]
\node[data]  (raw)    {Dados brutos\\FNN, UPFD, Bluesky};
\node[proc, right=of raw]   (graph)  {Constru\c{c}\~ao\\de grafos};
\node[proc, right=of graph] (feat)   {\emph{Features}\\posicionais\\+\,BERT};
\node[model, right=of feat] (model)  {GNN\\GCN/GAT/SAGE};
\node[eval, right=of model] (val)    {Valida\c{c}\~ao\\k-fold pareado};
\node[eval, right=of val]   (out)    {M\'etricas\\F1, Cohen's $d$};

\draw[arr] (raw)   -- (graph);
\draw[arr] (graph) -- (feat);
\draw[arr] (feat)  -- (model);
\draw[arr] (model) -- (val);
\draw[arr] (val)   -- (out);
\end{tikzpicture}}
\caption{Vis\~ao geral do \emph{pipeline} de detec\c{c}\~ao. Dados brutos s\~ao
transformados em grafos de propaga\c{c}\~ao com \emph{features} estruturais
(e textuais via BERT, quando dispon\'iveis), passados por modelos GNN, e
avaliados com protocolo de k-fold pareado.}
\label{fig:pipeline}
\end{figure}

% ------------------------------------------------------------
\section{Datasets}
\label{sec:datasets}

% ESCREVER (1 paragrafo introdutorio + tabela T0 abaixo).
% Citacoes: \cite{shu2020fakenewsnet}, \cite{dou2021upfd},
% \cite{kleppmann2024atproto}, \cite{failla2024bluesky}.

\begin{table}[H]\centering
\caption{Datasets utilizados, com tamanho, fonte e papel no trabalho.}
\label{tab:datasets}
\small
\begin{tabular}{lrrll}
\toprule
\textbf{Dataset} & \textbf{N grafos} & \textbf{N classes} & \textbf{Fonte} & \textbf{Papel} \\
\midrule
FakeNewsNet PolitiFact   & $\sim$754   & 2 (50/50)   & \cite{shu2020fakenewsnet}      & Treino + valida\c{c}\~ao \\
UPFD-PolitiFact          & $\sim$314   & 2 (50/50)   & \cite{dou2021upfd}             & Cross-dataset \\
UPFD-GossipCop           & $\sim$5{.}464 & 2 (50/50) & \cite{dou2021upfd}             & Cross-dataset \\
Bluesky (acad\^emico)    & 168{.}463   & --          & \cite{failla2024bluesky}       & Aplica\c{c}\~ao (sem \emph{labels}) \\
\bottomrule
\end{tabular}
\end{table}

% Para cada dataset, 1 paragrafo: o que e, quem coletou, quantos grafos,
% feature dominante. Para Bluesky: declarar explicitamente que NAO tem
% ground-truth fake/real e e usado em modo demonstracao no Cap 5.

% ------------------------------------------------------------
\section{Constru\c{c}\~ao de grafos: FakeNewsNet}
\label{sec:construcao_fnn}

A constru\c{c}\~ao dos grafos do FakeNewsNet \'e realizada pelo \emph{script}
\texttt{00\_construir\_grafos\_fakenewsnet.py}. O \emph{script} efetua o
\emph{download} direto dos CSVs do reposit\'orio oficial do KaiDMML
\cite{shu2020fakenewsnet}, gera o \emph{embedding} textual do t\'itulo de
cada not\'icia via \emph{Sentence-BERT} multil\'ingue (modelo
\texttt{paraphrase-multilingual-mpnet-base-v2}, dimens\~ao 768) e
\emph{cacheia} o resultado em disco para evitar recomputa\c{c}\~ao. Cada
not\'icia gera um grafo em estrela: o n\'o raiz \'e a not\'icia em si, e
seus filhos s\~ao os identificadores de \emph{tweets} listados pelo CSV
como tendo propagado a not\'icia.

A topologia em estrela (e n\~ao em \'arvore com profundidade) \'e
imposta pelo formato do CSV bruto, que fornece apenas uma lista plana de
\texttt{tweet\_ids} por not\'icia, sem rela\c{c}\~ao de parentesco entre
\emph{retweets} e sem \emph{timestamps} (caveat j\'a discutido em
\S\ref{sec:datasets}). Aplicam-se dois filtros: o m\'inimo de dois
\emph{tweets} de propaga\c{c}\~ao por not\'icia (descarta cascatas
degeneradas, sem o que o grafo de propaga\c{c}\~ao seria reduzido a um
\'unico n\'o) e o m\'aximo de 100 n\'os por grafo. A distribui\c{c}\~ao
bruta de tamanhos de cascata \'e fortemente assim\'etrica: a mediana
\'e de 30 \emph{tweets}, mas a m\'edia chega a 552 e o m\'aximo observado
a 29.060. O limite de 100 n\'os trunca, portanto, 50\% dos grafos do
\emph{dataset}; n\~ao se trata de remo\c{c}\~ao de \emph{outliers} raros,
mas de decis\~ao de projeto com efeito substancial. A justificativa \'e
estrutural: em um grafo-estrela, todos os n\'os-filho compartilham o mesmo
\emph{embedding} BERT (replicado a partir do t\'itulo), e folhas
adicionais al\'em de uma centena acrescentam informa\c{c}\~ao marginal
desprez\'ivel; al\'em disso, o \emph{global mean pool} adotado pela
camada final faria com que, em um grafo de dezenas de milhares de folhas
quase id\^enticas, a contribui\c{c}\~ao do n\'o raiz, \'unico portador
das \emph{features} \texttt{is\_root} e \texttt{grau\_norm}, fosse
dilu\'ida a uma fra\c{c}\~ao \'infima do vetor agregado. Como
contrapartida, os grafos truncados tornam-se indistingu\'iveis quanto ao
tamanho absoluto da cascata (limita\c{c}\~ao retomada em
\S\ref{sec:limitacoes}).

O \emph{script} oferece tr\^es variantes de constru\c{c}\~ao, controladas
pela \emph{flag} \texttt{--feature-variant}: \texttt{posfull} (todas as
tr\^es codifica\c{c}\~oes posicionais ativas, vetor nodal de 771
dimens\~oes), \texttt{posmin} (apenas \texttt{is\_root} e \texttt{pos}, com
\texttt{grau\_norm} zerado) e \texttt{posgrau} (apenas \texttt{is\_root} e
\texttt{grau\_norm}, com \texttt{pos} zerado). As tr\^es variantes
materializam-se em diret\'orios distintos em
\texttt{data/fakenewsnet\_<suffix>/} e s\~ao consumidas pela ablation
intra-encoding (\S\ref{sec:topo_sem_texto}). Uma asser\c{c}\~ao de sanidade
no construtor (\texttt{x.unique(dim=0).shape[0] >= min(num\_nos, 3)})
garante que os vetores nodais de um mesmo grafo n\~ao sejam id\^enticos,
prevenindo o cen\'ario de degenera\c{c}\~ao do \emph{message passing}
descrito em \S\ref{sec:degeneracao}.

% ------------------------------------------------------------
\section{Folds compartilhados}
\label{sec:folds}

O \emph{script} \texttt{gerar\_folds.py} instancia
\texttt{StratifiedKFold(n\_splits=10, shuffle=True, random\_state=42)} do
\emph{scikit-learn} sobre os 754 grafos do FakeNewsNet e persiste os
\'indices em \texttt{data/folds\_fnn.pt}. Esse arquivo \'e gerado
\emph{uma \'unica vez} e consumido por todos os experimentos do trabalho
que reportam compara\c{c}\~oes pareadas: o baseline textual
(\S\ref{sec:baseline_textual}), o diagn\'ostico de \emph{confound}
(\S\ref{sec:confound}), o teste de signific\^ancia entre arquiteturas
GNN, entre outros.

A motiva\c{c}\~ao \'e estat\'istica. O teste \emph{t} pareado
(\texttt{scipy.stats.ttest\_rel}) que usamos para comparar dois modelos
exige que cada par de m\'etricas comparado tenha sido obtido sobre os
\emph{mesmos exemplos}. Caso contr\'ario, a varia\c{c}\~ao entre folds se
mistura \`a diferen\c{c}a entre modelos, reduzindo o poder do teste e
tornando suas conclus\~oes question\'aveis. O compartilhamento de
\emph{folds} elimina essa fonte de ru\'ido e \'e pr\'e-requisito do
protocolo de avalia\c{c}\~ao formalizado em \S\ref{sec:protocolo}.

% ------------------------------------------------------------
\section{Arquiteturas GNN e hiperpar\^ametros}
\label{sec:arquiteturas}

Tr\^es arquiteturas representativas das principais fam\'ilias de Redes
Neurais de Grafos s\~ao avaliadas: a \emph{Graph Convolutional Network}
(GCN, \cite{kipf2017gcn}), com agrega\c{c}\~ao espectral normalizada pelo
grau; a \emph{Graph Attention Network} (GAT, \cite{velickovic2018gat}),
com pesos de aresta aprendidos via mecanismo de aten\c{c}\~ao; e o
GraphSAGE (\cite{hamilton2017graphsage}), com agrega\c{c}\~ao indutiva por
m\'edia dos vizinhos. Todas as tr\^es s\~ao implementadas no \emph{PyTorch
Geometric}~\cite{fey2019pyg} e compartilham a mesma arquitetura agregada:
tr\^es camadas convolucionais, com 64 canais ocultos por camada,
seguidas de \emph{global mean pool} sobre os n\'os, \emph{dropout} de 0,5
e um classificador linear final. Os demais hiperpar\^ametros, incluindo as poucas diferen\c{c}as entre as
arquiteturas, est\~ao na Tabela~\ref{tab:hyperparams}.

\textit{Para o UPFD, o vetor de \emph{features} de entrada de cada n\'o
\'e o vetor pr\'e-computado fornecido pelo dataset (variante
\texttt{profile}, com 10 dimens\~oes); para o FakeNewsNet, o vetor \'e o
de 771 dimens\~oes descrito em \S\ref{sec:construcao_fnn}; para os
experimentos de topologia pura da \S\ref{sec:topo_sem_texto}, ambos os
\emph{datasets} t\^em suas \emph{features} substitu\'idas por
\texttt{[is\_root, grau\_norm]} (dimens\~ao 2).}

\begin{table}[H]\centering
\caption{Hiperpar\^ametros usados nos modelos GNN. Defaults s\~ao iguais
entre arquiteturas; diferen\c{c}as anotadas na pr\'opria tabela.}
\label{tab:hyperparams}
\small
\begin{tabular}{lrrr}
\toprule
\textbf{Hiperpar\^ametro} & \textbf{GCN} & \textbf{GAT} & \textbf{SAGE} \\
\midrule
\texttt{hidden\_channels}     & 64     & 64     & 64 \\
\texttt{num\_layers}          & 3      & 3      & 3 \\
\texttt{dropout}              & 0{.}5  & 0{.}5  & 0{.}5 \\
\texttt{num\_heads} (atten\c{c}\~ao) & --     & 4      & -- \\
\texttt{pooling}              & \texttt{global\_mean\_pool} & idem & idem \\
\texttt{learning\_rate}       & 0{.}001 & 0{.}001 & 0{.}001 \\
\texttt{weight\_decay}        & $5\times10^{-4}$ & $5\times10^{-4}$ & $5\times10^{-4}$ \\
\texttt{batch\_size}          & 32     & 32     & 32 \\
\texttt{optimizer}            & Adam   & Adam   & Adam \\
\texttt{scheduler}            & \multicolumn{3}{c}{ReduceLROnPlateau (factor=0{.}5, patience=5)} \\
\texttt{early\_stopping}      & \multicolumn{3}{c}{patience=7 sobre F1-macro de valida\c{c}\~ao} \\
\texttt{max\_epochs}          & 30     & 30     & 30 \\
\texttt{seed}                 & 42 (com varredura para 10 \emph{seeds} em \S\ref{sec:topo_sem_texto}) \\
\bottomrule
\end{tabular}
\end{table}

A op\c{c}\~ao por hiperpar\^ametros id\^enticos entre as tr\^es arquiteturas
n\~ao busca otimizar cada uma ao m\'aximo (o que tornaria a compara\c{c}\~ao
dependente de quanto cada modelo foi ajustado), mas sim \emph{isolar o
efeito arquitetural}: se uma fam\'ilia se destaca consistentemente sob a
mesma configura\c{c}\~ao b\'asica, o sinal observado \'e atribu\'ivel \`a
arquitetura, n\~ao ao \emph{tuning}. Os defaults adotados s\~ao valores
convencionais na literatura de GNNs e suficientes para conduzir o
treinamento \`a converg\^encia dentro do m\'aximo de 30 \'epocas
estabelecido.

% ------------------------------------------------------------
\section{\emph{Baselines} n\~ao-GNN}
\label{sec:baselines}

Dois \emph{baselines} deliberadamente simples acompanham as arquiteturas
GNN ao longo do trabalho, cada um isolando uma dire\c{c}\~ao do espa\c{c}o
de hip\'oteses.

O primeiro \'e um classificador puramente \emph{textual}: a
\emph{Regress\~ao Log\'istica} aplicada sobre o \emph{embedding} BERT da
raiz da \'arvore de propaga\c{c}\~ao (isto \'e, o vetor de 768
dimens\~oes do t\'itulo da not\'icia). Esse modelo n\~ao tem acesso a
nenhuma informa\c{c}\~ao estrutural: ignora o grafo inteiro e v\^e apenas
o conte\'udo textual. Estabelece, portanto, o \emph{teto textual} do
problema; qualquer GNN que n\~ao supere significativamente este
\emph{baseline} n\~ao est\'a agregando valor sobre a informa\c{c}\~ao do
texto.

O segundo \'e um classificador puramente \emph{estrutural tabular}: um
\emph{Random Forest} com 200 \'arvores, treinado sobre apenas duas
\emph{features} agregadas por grafo: o n\'umero total de n\'os
(\texttt{num\_nodes}) e o grau do n\'o raiz (\texttt{grau\_root}). Esse
modelo encarna o \emph{confound trivial} da detec\c{c}\~ao de fake news em
redes: a observa\c{c}\~ao recorrente de que not\'icias falsas tendem a
viralizar mais, gerando cascatas maiores. Se um \emph{Random Forest}
alimentado apenas com tamanho e grau atinge um F1 pr\'oximo do da GNN, o
ganho aparente da GNN n\~ao vem de capturar estrutura sutil, mas apenas
de reproduzir esse \emph{confound}; se, por outro lado, o GNN supera
significativamente o RF, h\'a sinal estrutural genu\'ino sendo aprendido.
Os dois \emph{baselines} s\~ao confrontados, lado a lado, com as
arquiteturas GNN em \S\ref{sec:baseline_textual} e \S\ref{sec:confound}.

% ------------------------------------------------------------
\section{Protocolo de avalia\c{c}\~ao}
\label{sec:protocolo}

A m\'etrica prim\'aria adotada \'e o \emph{F1-macro}, a m\'edia n\~ao
ponderada dos F1 das duas classes (fake e real). Como os tr\^es
\emph{datasets} supervisionados utilizados s\~ao aproximadamente
balanceados em torno de $50\%/50\%$ entre as classes
(\S\ref{sec:datasets}), o F1-macro \'e numericamente pr\'oximo da
acur\'acia, mas com a vantagem de permanecer significativo caso futuros
\emph{datasets} apresentem desequil\'ibrio, situa\c{c}\~ao comum em
detec\c{c}\~ao de fake news em produ\c{c}\~ao, onde a classe
\emph{fake} pode ser minorit\'aria. A literatura UPFD~\cite{dou2021upfd}
reporta acur\'acia; usamos F1-macro para alinhar com pr\'aticas mais
recentes em classifica\c{c}\~ao bin\'aria desbalanceada e marcamos
expressamente a diferen\c{c}a quando confrontamos com valores publicados
(Tabela~\ref{tab:upfd_vs_publicado}).

Para comparar dois modelos sob o mesmo \emph{dataset}, utilizamos o
\emph{teste t pareado} (\texttt{scipy.stats.ttest\_rel}) sobre os
\emph{folds} compartilhados de \S\ref{sec:folds}. O pareamento tem
duas implica\c{c}\~oes pr\'aticas: a varia\c{c}\~ao entre \emph{folds}
\'e cancelada (cada \emph{fold} contribui com uma diferen\c{c}a entre
modelos, n\~ao com uma m\'etrica absoluta), e o teste recupera o sinal
mesmo em diferen\c{c}as pequenas que o teste n\~ao-pareado n\~ao
detectaria. No caso particular dos \emph{subsets} \emph{post-hoc} (como
os recortes por profundidade da cascata em \S\ref{sec:estratificacao}),
nos quais o pareamento por fold \'e imposs\'ivel por o \emph{subset} ser
definido sobre o conjunto de teste \'unico, recorremos ao \emph{bootstrap}
pareado descrito em \S\ref{sec:metod_estratificacao}.

Para a an\'alise descritiva de separabilidade estrutural entre fake e
real (\S\ref{sec:porque}), reportamos o \emph{tamanho de efeito de Cohen}
($d$), conforme defini\c{c}\~ao cl\'assica de
\cite{cohen1988statistical}. A escala interpretativa adotada \'e a usual:
$|d| < 0{,}2$ desprez\'ivel, $0{,}2 \leq |d| < 0{,}5$ pequeno,
$0{,}5 \leq |d| < 0{,}8$ m\'edio, $|d| \geq 0{,}8$ grande. A tese
central do trabalho est\'a operacionalizada nesse limiar (ver
\S\ref{sec:negativo}).

\textit{Caveat de reprodutibilidade.} Todos os experimentos foram
executados em CPU (Intel \emph{Raptor Lake}, 16 \emph{cores} l\'ogicos,
16~GB de RAM), sob Python~3.12.10 e PyTorch~2.10+cpu, com semente
universal 42. Os resultados s\~ao reproduz\'iveis no mesmo ambiente, mas
n\~ao bit-exato entre m\'aquinas: opera\c{c}\~oes \texttt{scatter} e
\texttt{gather} do PyTorch Geometric n\~ao possuem \emph{kernel}
determin\'istico em GPU, e diferen\c{c}as entre vers\~oes de CUDA/cuDNN
ou entre CPUs distintas geram varia\c{c}\~ao tipicamente na terceira casa
decimal das m\'etricas. Como contramedida, os achados centrais foram
validados em 10 \emph{seeds} independentes (\S\ref{sec:topo_sem_texto}) e
em \emph{bootstrap} pareado de 2.000 reamostragens
(\S\ref{sec:estrat_cascatas}); as conclus\~oes qualitativas se mant\^em
entre execu\c{c}\~oes.

% ------------------------------------------------------------
\section{An\'alise estratificada do erro}
\label{sec:metod_estratificacao}

O F1-macro agregado mede \emph{quanto} um modelo acerta, n\~ao \emph{onde}
ele acerta. Para responder a essa segunda quest\~ao, relevante para a
caracteriza\c{c}\~ao de regimes em que cada modelo falha ou prospera,
adotamos tr\^es t\'ecnicas complementares aplicadas sobre o
\emph{test set} do UPFD-GossipCop, em conjunto detalhadas em
\S\ref{sec:estratificacao}.

A primeira \'e uma \emph{regress\~ao localmente ponderada} (LOWESS) de
$P(\text{acerto} \mid m)$ em fun\c{c}\~ao de cada m\'etrica estrutural $m$.
Implementamo-la via \emph{rolling mean} de janela proporcional ($10\%$
da amostra ordenada, com m\'inimo absoluto de 30 observa\c{c}\~oes),
evitando a depend\^encia de \texttt{statsmodels}. A vantagem sobre o
\emph{binning} discreto \'e dupla: visualiza-se continuamente a varia\c{c}\~ao
da acur\'acia ao longo de $m$, sem decidir arbitrariamente onde cortar os
\emph{bins}, o que evita o espa\c{c}o para \emph{cherry-picking}.

A segunda \'e a defini\c{c}\~ao de \emph{outliers} via \emph{percentil
emp\'irico} (p5/p95) sobre o pr\'oprio \emph{test set}, em vez de
limiares baseados em \emph{z-score}. A escolha \'e justificada pela forma
das distribui\c{c}\~oes: m\'etricas de cascata em redes sociais s\~ao
\emph{heavy-tail}: a m\'edia n\~ao \'e robusta, o desvio-padr\~ao \'e
inflado por poucos casos extremos, e crit\'erios baseados em $\sigma$
tornam-se mal-definidos. O percentil emp\'irico \'e invariante \`a forma
da distribui\c{c}\~ao e identifica reciprocamente o $5\%$ inferior e o
$5\%$ superior independentemente da assimetria.

A terceira \'e o \emph{bootstrap pareado} (com $n_{\mathrm{boot}}=2.000$)
para intervalos de confian\c{c}a sobre a diferen\c{c}a de F1 entre dois
modelos avaliados no mesmo \emph{subset}, t\'ecnica empregada em
\S\ref{sec:estrat_cascatas} para confirmar a significancia do ganho do
SAGE estrutural sobre o RF tabular em recortes de profundidade da
cascata. O \emph{bootstrap} \'e o substituto natural do
\texttt{ttest\_rel} pareado por \emph{fold} quando o pareamento por
\emph{fold} n\~ao se aplica, como ocorre em \emph{subsets} definidos
\emph{post-hoc} sobre o teste \'unico do UPFD.

% ------------------------------------------------------------
\section{An\'alise de explicabilidade}
\label{sec:analise_explicabilidade}

Para inspecionar o que a arquitetura SAGE estrutural efetivamente aprende
do grafo de propaga\c{c}\~ao, recorremos ao \emph{GNNExplainer}
\cite{ying2019gnnexplainer}, implementado em
\texttt{torch\_geometric.explain.Explainer} com algoritmo
\texttt{GNNExplainer(epochs=200)} e \texttt{edge\_mask\_type="object"}, em
modo de classifica\c{c}\~ao multi-classe sobre o grafo inteiro. O
GNNExplainer aprende uma m\'ascara sobre as arestas do grafo que maximiza
a informa\c{c}\~ao m\'utua entre a sub-rede mascarada e a predi\c{c}\~ao do
modelo, identificando assim quais arestas s\~ao mais decisivas para a
classifica\c{c}\~ao de cada grafo.

A an\'alise \'e aplicada sobre 20 amostras do \emph{test set} do
UPFD-GossipCop estratificadas em quatro grupos de cinco grafos cada
(\emph{fake}-pequena, \emph{fake}-grande, \emph{real}-pequena,
\emph{real}-grande), permitindo identificar padr\~oes diferenciais por
classe e por tamanho. Para cada amostra, calcula-se a fra\c{c}\~ao de
massa de import\^ancia concentrada em arestas a $1$-hop (raiz $\to$ filhos
diretos), $2$-hops e $\geq 3$-hops da raiz. Os resultados, apresentados
em \S\ref{sec:explainer}, revelam uma assimetria sistem\'atica entre
\emph{fake} e \emph{real} no padr\~ao espacial de import\^ancia,
achado que sustenta interpreta\c{c}\~oes sobre o mecanismo de
classifica\c{c}\~ao aprendido pelo modelo.

% ------------------------------------------------------------
\section{Classificador dual: regra de combina\c{c}\~ao \emph{post-hoc}}
\label{sec:dual}

Para a aplica\c{c}\~ao pr\'atica documentada no
Ap\^endice~\ref{ap:aplicacao}, combinamos o classificador textual
(LogReg-BERT) e o classificador topol\'ogico (RF estrutural) por meio de
uma regra de combina\c{c}\~ao heur\'istica. Essa regra \textbf{n\~ao \'e uma decis\~ao metodol\'ogica a
priori}: trata-se de uma heur\'istica derivada \emph{post-hoc}, ou seja,
ap\'os a observa\c{c}\~ao emp\'irica dos resultados de
\S\ref{sec:concordancia}. A formula\c{c}\~ao exata da regra, os n\'umeros
que a justificam e as alternativas descartadas est\~ao naquela se\c{c}\~ao,
sob a equa\c{c}\~ao \eqref{eq:regra_dual}. A regra \'e mencionada aqui
porque a ferramenta web do Ap\^endice~\ref{ap:aplicacao} usa o
classificador combinado.

% ------------------------------------------------------------
\section{Reprodutibilidade}
\label{sec:reprodutibilidade}

A reprodutibilidade depende de fixar o ambiente, controlar as fontes de
aleatoriedade e expor as limita\c{c}\~oes ao determinismo. Primeiro,
o ambiente de execu\c{c}\~ao \'e fixado: Python~3.12.10, PyTorch~2.10+cpu,
PyTorch~Geometric~2.7.0, \emph{scikit-learn} e demais depend\^encias
versionadas em \texttt{requirements.txt} (\emph{pinned}). Os experimentos
foram conduzidos em \textbf{ambiente CPU} (Intel \emph{Raptor Lake}, 16
\emph{cores} l\'ogicos, 16~GB de RAM, Windows~11); n\~ao se utilizou
GPU em nenhum momento do treinamento ou da avalia\c{c}\~ao.

Segundo, todas as fontes de aleatoriedade s\~ao semeadas com o valor 42:
\texttt{torch.manual\_seed}, \texttt{numpy.random.seed},
\texttt{random.seed}, \texttt{StratifiedKFold(random\_state=42)} e os
construtores dos classificadores do \emph{scikit-learn} (\texttt{n\_jobs=-1}
\'e adotado para o RF, sem efeito sobre determinismo). Para os
experimentos sens\'iveis \`a inicializa\c{c}\~ao, conduzimos varredura
expl\'icita sobre 10 \emph{seeds} independentes em \S\ref{sec:topo_sem_texto}
e reportamos m\'edia e desvio-padr\~ao.

Terceiro, reconhecemos as limita\c{c}\~oes ao determinismo: operadores
\texttt{scatter} e \texttt{gather} do PyTorch Geometric n\~ao possuem
\emph{kernel} determin\'istico em GPU
\cite{ying2019gnnexplainer}\footnote{Refer\^encia da nota: a documenta\c{c}\~ao
oficial do PyTorch sobre reprodutibilidade
(\url{https://pytorch.org/docs/stable/notes/randomness.html}).}, e
diferen\c{c}as de arquitetura de CPU ou de vers\~ao de bibliotecas
geram varia\c{c}\~ao tipicamente na terceira casa decimal das m\'etricas
reportadas. As conclus\~oes qualitativas do trabalho s\~ao robustas a
essa varia\c{c}\~ao. O reposit\'orio oficial do c\'odigo, com o
\texttt{README.md} como fonte autoritativa de instru\c{c}\~oes, est\'a
publicado junto aos pesos persistidos em \cite{sibila2025bluesky}.